% ==================== NEW PAGE FOR PROPOSAL ====================
\newpage

% Remove top margin for this page only
\vspace*{-\topmargin}
\vspace*{-\headheight}
\vspace*{-\headsep}
\vspace*{-\topskip}

% ==================== PROPOSAL HEADER WITH FULL-WIDTH BACKGROUND ====================
\noindent
\colorbox{primary}{%
    \makebox[\linewidth][c]{%
        \parbox{\paperwidth}{%
            \vspace{0.8cm}
            
            % PROPOSAL heading
            \begin{center}
                {\Huge\bfseries\color{white} PROPOSAL}
            \end{center}
            
            \vspace{0.3cm}
            
            % Divider
            \begin{center}
                \textcolor{white}{\rule{0.6\textwidth}{0.5pt}}
            \end{center}
            
            \vspace{0.3cm}
            
            % Program and Project
            \begin{center}
                {\Large\bfseries\color{white} \program{} – Summer Term}
            \end{center}
            
            \vspace{0.2cm}
            
            \begin{center}
                {\large \color{white} \projecttitle}
            \end{center}
            
            \vspace{0.4cm}
            
            % Second divider
            \begin{center}
                \textcolor{white}{\rule{0.8\textwidth}{0.3pt}}
            \end{center}
            
            \vspace{0.3cm}
            
            % Project metadata
            \begin{center}
                {\color{white}\textbf{Project:}} \projecttitle
            \end{center}
            
            \begin{center}
                {\color{white}\textbf{Mentors:}} \mentors
            \end{center}
            
            \begin{center}
                {\color{white}\textbf{Duration:}} Mid-June to end of November 2026 (part-time, 12--15 hrs/week)
            \end{center}
            
            \vspace{0.3cm}
            
            % Third divider
            \begin{center}
                \textcolor{white}{\rule{0.8\textwidth}{0.3pt}}
            \end{center}
            
            \vspace{0.3cm}
            
            % Applicant info
            \begin{center}
                {\color{white}\textbf{\myname}}
            \end{center}
            
            \begin{center}
                {\color{white}\mydegree}
            \end{center}
            
            \begin{center}
                {\color{white}\href{\githuburl}{\color{white}\mygithub}}
            \end{center}
            
            \begin{center}
                {\color{white}\mylocation}
            \end{center}
            
            \begin{center}
                {\color{white}\texttt{\myemail}}
            \end{center}
            
            \vspace{0.8cm}
        }%
    }%
}

% Restore normal spacing after the full-width box
\vspace{0.5cm}

% ==================== INDEX / TABLE OF CONTENTS ====================
\begin{center}
    {\Large\bfseries\color{primary} Index / Table of Contents}
\end{center}

\vspace{0.2cm}
\begin{center}
    \textcolor{secondary}{\rule{0.5\textwidth}{0.3pt}}
\end{center}
\vspace{0.3cm}

\tableofcontents

\newpage
\section{Executive Summary}

Hyperledger Fabric-X is not a minor release upgrade of Hyperledger Fabric. It is a different network architecture with a different operational center of gravity. In classic Fabric, Fablo can describe a network in terms of organizations, peers, orderers, certificate authorities, channels, and chaincodes, then generate a working Docker-based environment from a single \texttt{fablo-config.json}. Fabric-X changes those assumptions. The Fabric-X repository and sample deployment material show a decomposed system with separate orderer, sidecar, coordinator, verifier, validator-committer, query service, PostgreSQL, and Fabric Smart Client nodes. The bundled \texttt{xdev} sample in \texttt{fabric-x-repo/samples/tokens/compose-xdev.yml} is already enough to prove that running Fabric-X locally requires more than renaming classic Fabric containers. Its startup path depends on service-specific environment variables, routing tables, FSC node configuration, and crypto material that is currently assembled from sample scripts rather than from a tool like Fablo.

Fablo is the right foundation for solving this because it already owns the part of the developer experience that matters most: taking a compact declarative configuration and turning it into reproducible local infrastructure. The current codebase uses a Docker-based generation approach where the \texttt{setup-network} command runs inside a container to generate network artifacts using EJS templates. However, this monolithic architecture makes it difficult to support fundamentally different network types like Fabric-X without extensive modifications.

The solution is to introduce a pluggable engine architecture that separates network-type-specific logic into interchangeable engines. This MVP will implement the engine framework and a Fabric-X engine that reads a single \texttt{fablo-config-fabricx.json}, validates it against a versioned JSON schema, generates all Docker Compose and configuration artifacts, and bootstraps a working local Fabric-X network. The engine architecture will maintain backward compatibility with classic Fabric while enabling clean extension for new network types.

\begin{tcolorbox}[colback=blue!5!white, colframe=blue!60!black,
  title=Core Deliverable, fonttitle=\bfseries, arc=4pt]
A pluggable engine architecture for Fablo with a production-quality Fabric-X engine that reads a single \texttt{fablo-config-fabricx.json}, validates it against a versioned JSON schema, generates all Docker Compose and configuration artifacts, and bootstraps a working local Fabric-X network with a single command --- with full lifecycle management (\texttt{fablo up}, \texttt{fablo down}, \texttt{fablo status}, \texttt{fablo reset}), automated e2e tests mirroring Fablo's existing test structure, and contributor documentation.
\end{tcolorbox}

The mentorship work will implement this engine architecture from scratch, starting with the current Docker-based Fablo codebase. The main themes are: introduce schema-first engine detection, isolate Fabric-X behavior in dedicated engine code, implement a clean local developer workflow for the bundled \texttt{xdev} topology, then extend toward richer topology, stronger crypto generation, and broader lifecycle coverage.

\begin{tcolorbox}[colback=gray!5!white, colframe=gray!60!black,
  title=Stretch Goals (begin Week 13 only after core is merged),
  fonttitle=\bfseries, arc=4pt]
Multi-org Fabric-X topologies; separate-container committer/endorser/orderer support; dynamic crypto generation pipeline.
\end{tcolorbox}
support; CI integration example; Fablo UI support for Fabric-X configs.
\end{tcolorbox}

\section{The Problem: Why This Work Matters}

\subsection{The Fabric-X developer experience gap today}

The current Fabric-X local developer workflow is fragmented across samples, scripts, and operational knowledge that a new contributor has to reconstruct manually. The repository explains the high-level architecture and points to supporting components, but it does not offer the kind of one-file, one-command bootstrap that Fablo users expect on the classic Fabric side. The sample topology under \texttt{fabric-x-repo/samples/tokens/compose-xdev.yml} uses a single image, \texttt{ghcr.io/hyperledger/fabric-x-committer-test-node:0.1.7}, but still requires the user to understand orderer identity variables such as \texttt{SC\_SIDECAR\_ORDERER\_IDENTITY\_MSP\_DIR}, channel selection via \texttt{SC\_SIDECAR\_ORDERER\_CHANNEL\_ID}, logging-level knobs for multiple internal components, and the expected crypto volume layout. The companion script \texttt{fabric-x-repo/samples/tokens/scripts/gen\_crypto.sh} shows that even the developer sample needs a sequence of Fabric CA bootstrap, node certificate enrollment, owner identity enrollment, and token parameter generation with \texttt{tokengen}. That is far beyond the onboarding budget of most evaluators or workshop participants.

The richer deployment path is even more demanding. The Ansible inventory in \texttt{fabric-x-repo/samples/tokens/ansible/inventory/fabric-x.yaml} is explicit about how much topology knowledge is hidden behind the current sample commands. It declares four separate orderer groups, each split into router, batcher, consenter, and assembler roles with distinct ports. It models committer services as separate hosts for the database, validator, verifier, coordinator, sidecar, and query service. It also contains a special \texttt{load\_generators} section used only to generate crypto material for users and genesis assembly. In other words, Fabric-X already has a real topology model, but that model is exposed to developers through infrastructure inventory, not through a simple application-facing configuration file.

This matters because developer tooling sets the adoption ceiling for an emerging platform. A contributor who wants to test a patch, run a workshop demo, or compare Fabric-X behavior with classic Fabric should not have to learn Ansible inventories, Fabric CA enrollment steps, and Token SDK crypto scripts before seeing a running network. The absence of a simple local workflow slows down issue reproduction, discourages documentation contributors, and makes project evaluation harder for exactly the people LFX mentorships are meant to empower.

\subsection{What Fablo does today and why it is the right foundation}

Fablo already solves the analogous problem for classic Hyperledger Fabric with a level of maturity visible directly in the repository. The configuration model is schema-driven and compact: \texttt{fablo-config.json} files contain \texttt{\$schema}, \texttt{global}, \texttt{orgs}, \texttt{channels}, \texttt{chaincodes}, and \texttt{hooks}. The types live in \texttt{src/types/FabloConfigJson.ts}. The classic generation path flows from \texttt{src/commands/generate/index.ts} into \texttt{src/engines/classic/index.ts}, which validates against the classic schema, writes a normalized config into the target directory, and invokes \texttt{node bin/run.mjs setup-network fablo-config.json}. That in turn reaches \texttt{src/setup-docker/index.ts}, whose \texttt{writing()} method expands the config, renders EJS templates, writes connection profiles, creates \texttt{fabric-docker/docker-compose.yaml}, writes scripts, and generates a Mermaid topology file.

The surrounding ecosystem reinforces that this is not a prototype tool. \texttt{fablo-repo/SUPPORTED\_FEATURES.md} lists tested support for BFT, RAFT, TLS, snapshots, CCaaS, connection profiles, channel query tooling, and a broad command surface including \texttt{generate}, \texttt{up}, \texttt{down}, \texttt{reset}, \texttt{validate}, \texttt{extend-config}, and \texttt{export-network-topology}. The e2e coverage is equally important. \texttt{fablo-repo/e2e-network/docker/test-01-v3-simple.sh} demonstrates the house style for integration tests: build first, initialize a config in a temp directory, bring the network up, wait for service readiness by checking container logs, run functional queries and invokes, exercise lifecycle commands, and tear the environment down. That pattern is exactly what a Fabric-X contribution should reuse rather than replace.

Fablo is therefore not just ``a CLI that starts Docker Compose''. It is a stable user contract, a test harness, and a generation pipeline with a clear place to insert specialized logic. The recent engine split makes that explicit. The Fabric-X work can now extend Fablo without contaminating the classic path.

\subsection{Why the integration is non-trivial}

The integration is non-trivial because the abstractions on the two sides are genuinely different. Classic Fabric configs describe organizations that own peers, certificate authorities, and orderers; channels bind organizations; chaincodes attach to channels; the generator emits the expected Docker assets for that world. Fabric-X does not center those concepts in the same way. The current schema in \texttt{src/engines/fabricx/schema/fabricx-schema-v1.json} asks for \texttt{channelId}, \texttt{namespace}, an infrastructure image, port assignments, and a list of FSC-style nodes with IDs, roles, API ports, and P2P ports. The generated node config in \texttt{src/engines/fabricx/templates/core.yaml.ejs} is oriented around \texttt{fsc}, routing configuration, identity resolvers, a Fabric-X driver, and token management. That is structurally different from classic Fabric peer definitions.

Lifecycle ordering is also different. The Fabric-X engine cannot safely mirror the classic ``generate on up'' habit because a decomposed runtime is more sensitive to partially generated state. The current implementation already made the correct move: \texttt{fablo up} for Fabric-X now requires explicit prior generation, checks that \texttt{docker-compose.xdev.yaml} and the engine state file exist, and then waits for the orderer and query ports to become reachable. Crypto handling is different again. Classic Fablo assumes the classic Fabric crypto toolchain. Fabric-X currently depends on pre-generated sample crypto, explicit orderer MSP environment variables, and a separate Token SDK namespace-generation flow. Those differences are not cosmetic. They are exactly why this project needs careful architecture rather than a thin wrapper around existing scripts.

\section{Technical Background and Research}

\subsection{Fablo codebase analysis}

\begin{table}[h]
\centering
\begin{tabularx}{\textwidth}{|l|X|}
\hline
\textbf{Layer} & \textbf{Finding from source reading} \\
\hline
Config schema & Classic configs use \texttt{\$schema}, \texttt{global}, \texttt{orgs}, \texttt{channels}, \texttt{chaincodes}, and \texttt{hooks} as defined in \texttt{src/types/FabloConfigJson.ts}. \\
\hline
Generation pipeline & Currently uses Docker-based generation: \texttt{setup-network} command runs in a container, calls \texttt{src/setup-docker/index.ts} which renders EJS templates from \texttt{src/setup-docker/templates/} to generate Docker Compose files, connection profiles, and scripts. \\
\hline
Command structure & Oclif-based CLI with commands like \texttt{setup-network}, \texttt{validate}, \texttt{init}, etc. Lifecycle operations (\texttt{up}, \texttt{down}, \texttt{start}, \texttt{stop}) are handled by generated shell scripts in \texttt{fablo-target/}. \\
\hline
Template system & Uses EJS templates in \texttt{src/setup-docker/templates/} for generating Docker Compose files, connection profiles, CA configs, and utility scripts. \\
\hline
E2E test pattern & \texttt{e2e-network/docker/test-01-v3-simple.sh} builds, initializes, runs, waits for logs, performs functional checks, exercises lifecycle commands, and tears down. \\
\hline
Type system & \texttt{src/types/FabloConfigJson.ts} defines the configuration structure with proper TypeScript interfaces. \\
\hline
\end{tabularx}
\caption{Current Fablo codebase --- confirmed architecture from source reading}
\label{tab:fablo-analysis}
\end{table}

The current code path from user input to running containers works as follows: A user calls \texttt{fablo up fablo-config.json}, which invokes the \texttt{fablo.sh} script. This script checks if network files exist in \texttt{fablo-target/}. If not, it runs the \texttt{setup-network} command inside a Docker container, mounting the config file and target directory. The \texttt{src/setup-docker/index.ts} command parses the config, extends it with \texttt{extendConfig()}, and generates all artifacts using EJS templates. Once generated, \texttt{fablo.sh} calls the generated \texttt{fablo-target/fabric-docker.sh up} script to start the network.

This Docker-based generation approach works well for classic Fabric but creates challenges for supporting different network types like Fabric-X. The monolithic \texttt{setup-docker/index.ts} contains logic specific to classic Fabric topology, making it difficult to extend for fundamentally different architectures without extensive conditional branching.

\subsection{Fabric-X architecture analysis}

\begin{table}[h]
\centering
\begin{tabularx}{\textwidth}{|l|X|}
\hline
\textbf{Component} & \textbf{Role and config surface} \\
\hline
Arma Orderer & The README describes an orderer stack made of routers, batchers, consenters, and assemblers. In the Ansible inventory, routers expose ports such as \texttt{7050}, \texttt{7150}, \texttt{7250}, and \texttt{7350}; batchers, consenters, and assemblers use adjacent ports. \\
\hline
Endorser & In the current MVP schema, endorser is an FSC node role. It receives a generated \texttt{core.yaml} with node identity, P2P endpoint, routing references, Fabric-X driver settings, and token namespace wiring. \\
\hline
Validator-Committer & The Ansible inventory models the validator as its own committer component with \texttt{committer\_component\_type: validator} and an RPC port such as \texttt{5100}. It represents the validation and commitment half of the decomposed commit path. \\
\hline
Sidecar & The sidecar is the Fabric-facing entry point for FSC nodes. The xdev sample configures it with orderer identity variables and a public port \texttt{4001}. The current engine points node \texttt{peers} entries at \texttt{localhost:\{sidecar\}}. \\
\hline
Query Service & The xdev sample exposes the query service at \texttt{7001} through \texttt{SC\_QUERY\_SERVICE\_SERVER\_ENDPOINT=:7001}. The current engine waits for this port on startup and places it under \texttt{queryService} in generated node config. \\
\hline
Coordinator & The Ansible inventory models a separate \texttt{coordinator} component with its own RPC port, such as \texttt{5300}, and the xdev sample exposes a dedicated logging level variable for it. \\
\hline
Verifier & The verifier is also a distinct committer component in the inventory, typically with its own RPC port, such as \texttt{5200}. The xdev sample exposes \texttt{SC\_VERIFIER\_LOGGING\_LEVEL}. \\
\hline
PostgreSQL & The bundled xdev sample exposes PostgreSQL on \texttt{5433}. The decomposed inventory models it separately as \texttt{committer-db} with user, password, database name, and a different port (\texttt{5435}) in that topology. \\
\hline
\end{tabularx}
\caption{Fabric-X decomposed components --- confirmed from fabric-x-repo}
\label{tab:fabricx-components}
\end{table}

My working MVP deliberately targets the bundled \texttt{xdev} topology because that is the smallest viable integration surface that still represents real Fabric-X behavior. The sample image in \texttt{compose-xdev.yml} already bundles database, orderer, and committer behavior behind one container command, \texttt{run db orderer committer}. That is ideal for a first merge because it keeps the number of containers low while still requiring real Fabric-X configuration. The current Fablo engine maps onto this model directly: one generated compose file for the bundled infrastructure service and one \texttt{core.yaml} per logical node. This is enough to prove engine architecture, lifecycle delegation, template rendering, default port handling, and schema-first detection without taking on the full burden of a decomposed multi-service topology on day one.

At the same time, the local repository makes it clear that \texttt{xdev} is only a starting point. The Ansible inventory proves that Fabric-X's natural topology is decomposed, multi-process, and potentially multi-organization. The \texttt{tools/fxconfig} tree and its integration tests (\texttt{single\_org\_test.go} and \texttt{multi\_org\_test.go}) also show that configuration assembly and namespace lifecycle are already serious enough to justify dedicated tooling. That is why I should frame the MVP honestly: it is a local bootstrap path for the bundled developer topology, not the full end-state of Fabric-X operations.

\begin{figure}[h]
\centering
\fbox{\parbox{0.88\textwidth}{\centering\vspace{2.5cm}
\textit{[DIAGRAM PLACEHOLDER: Classic Fabric (peer + orderer + CA per org)
vs.\ Fabric-X decomposed architecture (Arma orderer, endorser, validator-committer,
sidecar, query service, coordinator, verifier)]}
\vspace{2.5cm}}}
\caption{Classic Hyperledger Fabric vs.\ Fabric-X decomposed architecture}
\label{fig:arch-comparison}
\end{figure}

\section{Proposed Solution: Implementing Pluggable Engine Architecture for Fabric-X}

\subsection{Integration path evaluation and decision}

\begin{table}[h]
\centering
\begin{tabularx}{\textwidth}{|p{2.6cm}|X|X|p{1.6cm}|}
\hline
\textbf{Approach} & \textbf{Advantages} & \textbf{Disadvantages} & \textbf{Verdict} \\
\hline
Separate repository & Total isolation from classic Fabric code; freedom to tailor commands, schemas, and templates only for Fabric-X; lower short-term coupling to historic Fablo assumptions. & Duplicates Fablo's CLI contract, target-directory handling, config parsing, docs surface, and test style; splits the contributor community; turns Fabric-X support into a parallel tool instead of an extension of the standard Fabric bootstrap workflow. & Rejected \\
\hline
Wrapper over Ansible & Fastest route to a demo if the only goal is to invoke existing sample tooling; preserves upstream sample scripts without deep integration work. & Leaves the developer with a two-layer mental model; errors surface through wrapper indirection; does not resolve schema mismatch; keeps ownership of generated artifacts outside Fablo; fails to create a real pluggable architecture. & Rejected \\
\hline
Pluggable engine in Fablo & Reuses mature commands, test patterns, target-dir logic, and documentation expectations; keeps classic behavior intact; isolates Fabric-X logic under \texttt{src/engines/fabricx/}; allows a future third engine without changing core CLI commands. & Requires implementing the engine architecture from scratch; needs careful handling of the existing Docker-based generation approach. & \textbf{Selected} \\
\hline
\end{tabularx}
\caption{Integration approach evaluation}
\label{tab:approaches}
\end{table}

The pluggable engine is the correct architectural choice because it enables clean support for different network types while maintaining Fablo's established user experience. The current Docker-based generation approach works for classic Fabric but cannot easily accommodate Fabric-X's different topology and configuration requirements. By introducing a pluggable engine architecture, we can isolate network-type-specific logic while preserving the existing CLI interface and generation workflow.

This approach will implement the engine framework as part of the MVP, creating a foundation that supports both classic Fabric and Fabric-X. The engine architecture follows the Open/Closed Principle: the core command code remains closed for modification with respect to network type, but open for extension through new engine implementations.

\subsection{Schema-first engine detection}

The key design choice is to detect the engine from the configuration schema, not from a custom \texttt{platform} field. This approach eliminates ambiguous states where a config might be misinterpreted. The schema identifies both the config shape and the engine that must process it.

\begin{tcolorbox}[colback=gray!5!white, colframe=gray!50!black,
  title=Engine detection logic (\texttt{src/engines/engine-loader.ts}),
  fonttitle=\bfseries, arc=4pt]
\begin{lstlisting}[language=javascript, basicstyle=\ttfamily\small,
  keywordstyle=\color{blue}, stringstyle=\color{orange}]
export function getEngine(config: any): FabloEngine {
  if (config === null || typeof config !== "object") {
    throw new Error("Invalid config: expected a JSON object.");
  }

  if (typeof config.$schema === "string" &&
      config.$schema.endsWith("fabricx-schema-v1.json")) {
    return new FabricXEngine();
  }
  return new ClassicFabricEngine();
}
\end{lstlisting}
\end{tcolorbox}

The suffix match is deliberate. Using \texttt{.endsWith()} instead of \texttt{.includes()} avoids false positives where an unrelated URL merely contains the Fabric-X schema name as a fragment or query component. Including \texttt{v1} in the schema filename creates a stable forward-compatibility story for future schema versions.

This schema-first detection will be implemented as part of the MVP, transitioning from the current monolithic Docker-based generation to a pluggable engine system. The new approach isolates engine selection to the schema itself, allowing the CLI commands to remain engine-agnostic while enabling clean code ownership for different network types.

\subsection{The FabloEngine interface}

The power of the engine approach comes from the fact that the interface is intentionally small and lifecycle-shaped rather than generation-shaped. This design matters because different network types have different startup requirements and lifecycle semantics.

\begin{tcolorbox}[colback=gray!5!white, colframe=gray!50!black,
  title=Engine contract (\texttt{src/engines/engine.ts}),
  fonttitle=\bfseries, arc=4pt]
\begin{lstlisting}[language=javascript, basicstyle=\ttfamily\small]
export interface ValidationResult {
  valid: boolean;
  errors: string[];
}

export interface FabloEngine {
  validate(rawConfig: unknown): ValidationResult;
  generate(config: any, targetDir?: string): Promise<void>;
  up(targetDir?: string): Promise<void>;
  down(targetDir?: string): Promise<void>;
  status(targetDir?: string): Promise<string>;
}
\end{lstlisting}
\end{tcolorbox}

\texttt{validate()} gives each engine control over its own schema and error formatting. \texttt{generate()} owns the conversion from a declarative config to runnable artifacts. \texttt{up()}, \texttt{down()}, and \texttt{status()} cover the minimum lifecycle surface that both engines must support. For classic Fabric, these will delegate to the existing generated \texttt{fabric-docker.sh} scripts. For Fabric-X, they will call Docker Compose directly and check for generated files first.

The FabricXEngine will use AJV for JSON schema validation against \texttt{fabricx-schema-v1.json}, ensuring type safety and error reporting. Templates will be rendered using EJS, with context variables passed directly from the validated config.

\subsection{The config schema design}

The current Fabric-X schema is intentionally narrow. It aims to make the bundled local topology ergonomic first, while reserving room for later decomposition. Every field in the sample config below comes from files actually present in this workspace: the \texttt{\$schema} constant from \texttt{fabricx-schema-v1.json}, the \texttt{global.tls} flag from the current engine and templates, the xdev image lineage from \texttt{compose-xdev.yml}, and the node role set from the TypeScript types and generated \texttt{core.yaml} template.

\begin{tcolorbox}[colback=gray!5!white, colframe=gray!50!black,
  title=Fabric-X network config (\texttt{fablo-config-fabricx.json}),
  fonttitle=\bfseries, arc=4pt]
\begin{lstlisting}[language=json, basicstyle=\ttfamily\footnotesize]
{
  "$schema": "https://fablo.io/schemas/fabricx-schema-v1.json",
  "global": {
    "fabricVersion": "x",
    "tls": false
  },
  "fabricx": {
    "channelId": "mychannel",
    "namespace": "token_namespace",
    "infrastructure": {
      "image": "ghcr.io/hyperledger/fabric-x-committer-test-node:0.1.7",
      "ports": {
        "orderer": 7050,
        "sidecar": 4001,
        "query": 7001,
        "database": 5433
      }
    },
    "nodes": [
      { "id": "issuer1", "type": "issuer", "apiPort": 9200, "p2pPort": 9201 },
      { "id": "endorser1", "type": "endorser", "apiPort": 9300, "p2pPort": 9301 },
      { "id": "owner1", "type": "owner", "apiPort": 9400, "p2pPort": 9401 }
    ]
  }
}
\end{lstlisting}
\end{tcolorbox}

This design is intentionally honest about what it is not yet modeling. There is no attempt to flatten decomposed orderer groups, validator services, or namespace generation settings into the MVP schema prematurely. Those belong in the next phase of the mentorship, where I will extend the config model for richer topology, stronger identity handling, and broader operational settings. The result is a schema that is small enough to teach, but structurally aligned with Fabric-X from the start.

The schema enforces required fields like \texttt{channelId}, \texttt{namespace}, and \texttt{nodes}, with optional \texttt{wallets} for each node to support identity management.

\subsection{Generated artifacts and docker-compose structure}

The MVP does not generate a cluster of separate Fabric-X service containers yet. Instead, it generates the smallest working set of artifacts needed for the bundled \texttt{xdev} path. That includes a Docker Compose file for the infrastructure container, a routing table, one node configuration per declared logical node (including specific \texttt{core.yaml} files for Fabric-X FSC nodes and integrating orderer configuration into the Docker Compose setup), copied or future-generated crypto material, and an engine state file used by lifecycle commands. This is a sound first merge because it proves the end-to-end workflow without pretending the bundled topology is the same as the eventual decomposed one.

\begin{tcolorbox}[colback=gray!5!white, colframe=gray!50!black,
  title=Generated \texttt{docker-compose.xdev.yaml} (excerpt),
  fonttitle=\bfseries, arc=4pt]
\begin{lstlisting}[language=yaml, basicstyle=\ttfamily\footnotesize]
services:
  committer:
    image: ghcr.io/hyperledger/fabric-x-committer-test-node:0.1.7
    container_name: committer
    networks:
      - fabricx
    environment:
      - SC_SIDECAR_ORDERER_IDENTITY_MSP_DIR=/root/config/crypto/peerOrganizations/org1.example.com/users/Admin@org1.example.com/msp
      - SC_SIDECAR_ORDERER_IDENTITY_MSP_ID=Org1MSP
      - SC_SIDECAR_ORDERER_CHANNEL_ID=mychannel
      - SC_SIDECAR_ORDERER_SIGNED_ENVELOPES=true
      - SC_QUERY_SERVICE_SERVER_ENDPOINT=:7001
      - SC_ORDERER_BLOCK_SIZE=1
    ports:
      - "4001:4001"
      - "7001:7001"
      - "7050:7050"
      - "5433:5433"
    volumes:
      - ./crypto:/root/config/crypto
    command: run db orderer committer
\end{lstlisting}
\end{tcolorbox}

This compose excerpt is complemented by \texttt{routing-config.yaml}, which enumerates node IDs and API endpoints, and by per-node \texttt{core.yaml} files. The generated node configuration wires FSC identity, websocket P2P addresses, sqlite persistence, endpoint resolvers for every peer node, Fabric-X driver selection, sidecar and query addresses, and token-management settings such as channel and namespace. The result is enough for \texttt{fablo generate} to materialize a coherent local developer environment from one JSON file.

\begin{figure}[h]
\centering
\fbox{\parbox{0.88\textwidth}{\centering\vspace{2.5cm}
\textit{[DIAGRAM PLACEHOLDER: fablo-config-fabricx.json
$\rightarrow$ FabricXEngine.generate()
$\rightarrow$ docker-compose.xdev.yaml + routing-config.yaml + per-node core.yaml
$\rightarrow$ fablo up $\rightarrow$ running Fabric-X network with port-readiness polling]}
\vspace{2.5cm}}}
\caption{Fablo-FabricX integration: config to running network}
\label{fig:proposed-arch}
\end{figure}

\section{MVP Plan and Timeline}

\subsection{Phase plan and milestone table}

\begin{table}[h]
\centering
\begin{tabularx}{\textwidth}{|c|p{3.2cm}|X|}
\hline
\textbf{Dates} & \textbf{Phase} & \textbf{Deliverables} \\
\hline
June 01--June 14 & Onboarding & Finalize the design proposal upstream, set up the local development environment, confirm PR review expectations with mentors, and rehearse the existing build and e2e harness so new Fabric-X work lands in the same workflow as classic Fablo changes. \\
\hline
June 15--July 31 & Phase 1: Implement Pluggable Engine Architecture & Create the \texttt{FabloEngine} interface and engine loader system; implement schema-first engine detection; develop the \texttt{ClassicFabricEngine} wrapper for existing Docker-based generation; create the \texttt{FabricXEngine} with AJV validation against \texttt{fabricx-schema-v1.json}; implement EJS templates for Fabric-X Docker Compose and configuration files; add \texttt{fablo init --fabricx} command; ensure backward compatibility with existing classic Fabric workflows. \\
\hline
August 01--August 23 & Phase 2: Lifecycle and Local Bootstrap & Wire up \texttt{fablo up}, \texttt{fablo down}, and \texttt{fablo status} against generated Fabric-X artifacts; finalize standard port polling, TLS handling, and precondition checks; and harden \texttt{e2e-network/fabricx/test-01-fabricx-simple.sh} so it validates the Docker Compose stack starts and tears down correctly, mirroring the existing \texttt{test-01-simple.sh} style E2E testing patterns. \\
\hline
August 24--August 31 & Midterm Evaluation & Deliver a working, mergeable PR that demonstrates the pluggable engine architecture with Fabric-X support for the bundled \texttt{xdev} topology, including schema validation, generated Docker assets, lifecycle commands, and automated e2e tests. \\
\hline
September 01--October 15 & Phase 3: The Crypto Pipeline and Routing & Replace copied sample crypto with dynamic generation orchestrated from Docker utility containers, generate token namespace parameters with \texttt{tokengen}, formalize orderer identity inputs and outputs, and deepen the generated routing layer so node and sidecar connectivity are driven directly from the Fabric-X config model. \\
\hline
October 16--November 14 & Phase 4: Multi-Service Topology and Hardening & Introduce \texttt{topology: decomposed} support for separate orderer, endorser, committer-family, and query-service containers; add stricter reset semantics and CCaaS-aware cleanup; and finish contributor documentation, user guides, and architectural decision records. \\
\hline
November 15--November 30 & Final Evaluation & Reserve buffer for final upstream review and integration work. Only if the prior phases are fully merged, use this period for stretch goals such as CI integration examples or Fablo UI exposure for Fabric-X networks. \\
\hline
\end{tabularx}
\caption{Official LFX timeline mapped to project milestones}
\label{tab:timeline}
\end{table}

This schedule reflects the need to implement the pluggable engine architecture from scratch while building on Fablo's existing foundation. Phase 1 focuses on creating the engine framework and the initial Fabric-X engine implementation. The midterm evaluation provides a concrete demonstration of the architecture working end-to-end. Subsequent phases extend the crypto generation, topology support, and hardening.

\subsection{Gantt placeholder}

\begin{figure}[h]
\centering
\fbox{\parbox{0.88\textwidth}{\centering\vspace{2.5cm}
\textit{[DIAGRAM PLACEHOLDER: Gantt chart --- June through November timeline
showing phases, milestones, and dependencies]}
\vspace{2.5cm}}}
\caption{Project Gantt chart}
\label{fig:gantt}
\end{figure}

The dependency structure behind the timeline is simple and defensible. Schema design depends on settled engine boundaries. Reliable lifecycle commands depend on stable generated artifacts. Docker-orchestrated crypto generation depends on understanding the existing Fabric-X sample toolchain. Decomposed topology support depends on first proving the engine model against \texttt{xdev}. Documentation and community adoption depend on all of the above.

\section{Full Technical Approach: Beyond the MVP}

\subsection{Crypto material generation for Fabric-X}

The MVP copies crypto material from sample locations because that is the smallest path to proving the engine contract. The full implementation should replace that with a Fabric-X-aware crypto pipeline built into the Fabric-X engine. The local sample script \texttt{fabric-x-repo/samples/tokens/scripts/gen\_crypto.sh} provides the roadmap. It starts a Fabric CA server, enrolls a CA admin, registers and enrolls owner identities using Idemix, enrolls TLS identities for FSC nodes, and finally runs \texttt{go tool tokengen gen zkatdlognogh.v1} to generate issuer and namespace parameters. That script is not optional sample decoration; it is the clearest current specification of how Fabric-X local identities are assembled.

The critical UX constraint is that this pipeline cannot depend on host-installed Go or Fabric-X binaries. Fablo's promise is that the user needs Docker and Node.js, not a handcrafted local toolchain. The right implementation is therefore Docker-orchestrated crypto generation. \texttt{FabricXEngine.generate()} should spin up ephemeral utility containers that mount the target directory, run the Fabric-X crypto scripts and binaries inside the container image, write artifacts into mounted volumes, and then exit. The user sees one \texttt{fablo generate} command; inside that command the engine runs one-off containers for \texttt{cryptogen}, Fabric CA enrollment, and \texttt{tokengen} steps.

I break this into three stages, and that is still the right structure. Stage one generates classic Fabric MSP and TLS material for orderer-side compatibility inside a utility container. Stage two generates FSC node identities and owner identities using Fabric CA enrollment logic derived from \texttt{gen\_crypto.sh}, again inside a disposable container or container pair. Stage three generates token namespace parameters with \texttt{tokengen} in a dedicated tool container. Each stage should have explicit output directories and validation checks so the engine can fail fast when a required artifact is missing.

\begin{tcolorbox}[colback=gray!5!white, colframe=gray!50!black,
  title=Planned crypto generation contract,
  fonttitle=\bfseries, arc=4pt]
\begin{lstlisting}[language=javascript, basicstyle=\ttfamily\small]
export interface FabricXCryptoOutputs {
  ordererMspDir: string;
  ordererMspId: string;
  nodeCertDirs: Record<string, string>;
  ownerWalletDirs: Record<string, string>;
  namespaceDir: string;
  genesisBlockPath?: string;
}

export interface FabricXCryptoGenerator {
  generate(config: FabricXConfig, targetDir: string): Promise<FabricXCryptoOutputs>;
  validate(outputs: FabricXCryptoOutputs): Promise<void>;
}
\end{lstlisting}
\end{tcolorbox}

This keeps the complexity isolated. The command layer still sees only \texttt{generate()}, but inside the Fabric-X engine the generator evolves from ``copy known crypto'' to ``run disposable tool containers and validate required identities''. That path is also consistent with the local Fabric-X version story. The current sample image is pinned to \texttt{0.1.7}, while \texttt{fabric-x-repo/go.mod} already depends on a newer committer lineage. Making crypto generation an explicit Docker-orchestrated module rather than a pile of copied files will make that version drift easier to manage.

\subsection{Multi-service topology support}

The current schema is effectively \texttt{xdev}-only. The Ansible inventory makes clear that a realistic Fabric-X topology includes separate orderer groups and separate committer family services. The next schema revision should therefore add an explicit topology mode and a structured infrastructure section for decomposed services. The important design point is that this should extend the existing Fabric-X section rather than reintroduce classic Fabric abstractions under different names.

The routing layer is where this becomes operationally real. Today, \texttt{routing-config.yaml.ejs} iterates over \texttt{nodes} and emits one route per FSC node. In the decomposed model, the generator has to emit both FSC-node routes and service-family endpoints so that sidecars and application-facing nodes can resolve the right infrastructure addresses for endorsers, committers, and query services. In practice, this means the EJS templates will map user-declared nodes such as \texttt{endorser1} or \texttt{owner1} together with decomposed service definitions into a deterministic routing table. The config is the source of truth; the generated routing file is a derived artifact.

\begin{tcolorbox}[colback=gray!5!white, colframe=gray!50!black,
  title=Planned multi-service Fabric-X config (illustrative),
  fonttitle=\bfseries, arc=4pt]
\begin{lstlisting}[language=json, basicstyle=\ttfamily\footnotesize]
{
  "$schema": "https://fablo.io/schemas/fabricx-schema-v2.json",
  "global": { "fabricVersion": "x", "tls": true },
  "fabricx": {
    "topology": "decomposed",
    "channelId": "arma",
    "namespace": "token_namespace",
    "ordererOrgs": [
      {
        "name": "OrdererOrg1",
        "domain": "ordererorg1.example.com",
        "group": 1,
        "routerPort": 7050,
        "batcherPort": 7051,
        "consenterPort": 7052,
        "assemblerPort": 7053
      }
    ],
    "committer": {
      "databasePort": 5435,
      "validatorPort": 5100,
      "verifierPort": 5200,
      "coordinatorPort": 5300,
      "sidecarPort": 4001,
      "queryPort": 7001
    },
    "nodes": [
      { "id": "issuer1", "type": "issuer", "apiPort": 9200, "p2pPort": 9201 },
      { "id": "endorser1", "type": "endorser", "apiPort": 9300, "p2pPort": 9301 },
      { "id": "owner1", "type": "owner", "apiPort": 9400, "p2pPort": 9401 }
    ]
  }
}
\end{lstlisting}
\end{tcolorbox}

I should not try to ship all of this at once. The correct sequencing is to land the \texttt{xdev} engine first, then add decomposed Docker Compose support for one topology slice, most likely the committer family, before tackling multi-organization orderer groups. This is both technically safer and easier to review upstream.

\subsection{CCaaS and endorser integration}

Classic Fablo already treats chaincode lifecycle as a first-class concern, and Fabric-X needs an equivalent story rather than a silent omission. The relevant difference is that Fabric-X's execution path is centered on the endorser microservice and Chaincode-as-a-Service style connectivity instead of the classic peer lifecycle. That means the Fabric-X config model eventually needs a dedicated section for contract execution inputs, not just infrastructure and node topology.

The practical MVP extension is to let the Fabric-X config declare one or more CCaaS contracts with a local source directory, image or build instructions, and the endpoint that the endorser-side execution environment should use. Fablo should then generate the necessary Docker Compose services or mounts so a developer can bind a local chaincode directory into the Fabric-X execution environment in the same spirit as the existing classic dev-mode and CCaaS support. The goal is to preserve the Fablo user expectation that a local project directory can be part of the declarative network description.

\begin{tcolorbox}[colback=gray!5!white, colframe=gray!50!black,
  title=Planned CCaaS extension for Fabric-X,
  fonttitle=\bfseries, arc=4pt]
\begin{lstlisting}[language=json, basicstyle=\ttfamily\footnotesize]
{
  "fabricx": {
    "chaincodes": [
      {
        "name": "tokencc",
        "lang": "ccaas",
        "directory": "./chaincodes/tokencc",
        "endpoint": "tokencc:9999",
        "endorser": "endorser1"
      }
    ]
  }
}
\end{lstlisting}
\end{tcolorbox}

This is not only a convenience feature. It is required if the Fabric-X path is supposed to support realistic developer testing. A local network that can boot but cannot attach and exercise a contract is incomplete from the perspective of Fablo users.

\subsection{Full e2e test suite}

The Fabric-X test strategy should consciously mimic the classic Fablo style rather than inventing a new harness. The repository already shows the pattern in \texttt{e2e-network/docker/test-01-v3-simple.sh}: create a temp workspace, build, initialize, start, wait, assert, mutate, and tear down. The existing Fabric-X test in \texttt{e2e-network/fabricx/test-01-fabricx-simple.sh} is already aligned with that spirit, and the mentorship should grow it into a small matrix of focused scenarios: generation-only validation, happy-path xdev lifecycle, missing-crypto behavior, config validation failures, and eventually decomposed topology bootstrap.

\begin{tcolorbox}[colback=gray!5!white, colframe=gray!50!black,
  title=E2E test: \texttt{test-01-fabricx-simple.sh} (full structure),
  fonttitle=\bfseries, arc=4pt]
\begin{lstlisting}[language=bash, basicstyle=\ttfamily\footnotesize]
#!/usr/bin/env bash
# test-01-fabricx-simple.sh
# Validates: single-org xdev Fabric-X network bootstraps and responds

set -euo pipefail

ROOT_DIR="$(cd "$(dirname "${BASH_SOURCE[0]}")/../.." && pwd)"
WORKDIR="$ROOT_DIR/e2e-network/__tmp__/fabricx-test-01"

rm -rf "$WORKDIR"
mkdir -p "$WORKDIR"
cd "$WORKDIR"

echo "== Build =="
(cd "$ROOT_DIR" && npm run build >/dev/null)

echo "== Init =="
npx fablo init --fabricx
test -f fablo-config-fabricx.json
grep -q 'fabricx-schema-v1.json' fablo-config-fabricx.json

CONFIG_FILE="$WORKDIR/fablo-config-fabricx.json"
TARGET_DIR="$WORKDIR/fablo-target/fabricx"

echo "== Generate =="
npx fablo generate "$CONFIG_FILE" "$TARGET_DIR"
test -f "$TARGET_DIR/docker-compose.xdev.yaml"
test -f "$TARGET_DIR/conf/routing-config.yaml"
test -f "$TARGET_DIR/conf/issuer1/core.yaml"

echo "== Up =="
FABLO_LOCAL=true npx fablo up "$CONFIG_FILE" "$TARGET_DIR"

echo "== Wait For Ports =="
for port in 7050 7001; do
  for i in $(seq 1 60); do
    nc -z 127.0.0.1 "$port" && break
    sleep 1
  done
  nc -z 127.0.0.1 "$port" || exit 1
done

echo "== Status =="
npx fablo status "$CONFIG_FILE" "$TARGET_DIR" | grep -q 'committer'

echo "== Down =="
npx fablo down "$CONFIG_FILE" "$TARGET_DIR"
docker ps --filter name=committer --format '{{.Names}}' | grep -q '^committer$' && exit 1

echo "ALL TESTS PASSED"
\end{lstlisting}
\end{tcolorbox}

This structure is intentionally simple. It proves the exact user contract the mentorship cares about. More ambitious functional token-flow checks can be added later, but the first responsibility of the test suite is to guarantee that Fablo can reliably bootstrap and tear down a Fabric-X network without regressing classic Fabric behavior.

\subsection{Lifecycle management: reset and channel tooling}

The current interface does not yet include a Fabric-X-specific \texttt{reset()} method, but the project scope should. Classic Fabric reset is largely about containers, chaincode state, and generated network artifacts. Fabric-X reset has to reason about a different state surface: Docker volumes, the bundled database state, generated per-node configuration, generated crypto material, generated namespace artifacts, and possible CCaaS-related containers or mounts. A bare \texttt{docker compose down -v} is not strict enough on its own.

\begin{tcolorbox}[colback=gray!5!white, colframe=gray!50!black,
  title=Planned Fabric-X reset flow,
  fonttitle=\bfseries, arc=4pt]
\begin{lstlisting}[language=javascript, basicstyle=\ttfamily\small]
async reset(targetDir?: string): Promise<void> {
  const dir = resolveTargetDir(targetDir);
  const composeFile = path.join(dir, "docker-compose.xdev.yaml");

  if (await fs.pathExists(composeFile)) {
    execSync(`docker compose -f ${composeFile} down -v`, { stdio: "inherit" });
  }

  execSync(`docker container prune -f`, { stdio: "inherit" });
  await fs.remove(path.join(dir, "conf"));
  await fs.remove(path.join(dir, "crypto"));
  await fs.remove(path.join(dir, "data"));
  await fs.remove(path.join(dir, "namespace"));
  await fs.remove(path.join(dir, "wallets"));
  await fs.remove(path.join(dir, "fabricx-engine-state.json"));
}
\end{lstlisting}
\end{tcolorbox}

The cleanup contract should be explicit in the documentation: \texttt{fablo reset} for Fabric-X tears the stack down, removes generated \texttt{conf/}, \texttt{crypto/}, \texttt{data/}, \texttt{namespace/}, and wallet state under the target directory, and cleans up container state associated with the generated network. Channel tooling is a related future concern. Classic Fablo already exposes channel query helpers because classic Fabric has a channel administration surface worth scripting. Fabric-X has its own namespace and configuration lifecycle through \texttt{fxconfig}. The right medium-term direction is not to shoehorn classic channel tooling into Fabric-X, but to add Fabric-X-aware helper commands once the bootstrap engine is stable.

\section{Known Limitations and How They Will Be Resolved}

\begin{tcolorbox}[colback=yellow!5!white, colframe=yellow!60!black,
  title=Known MVP Limitations (all isolated inside \texttt{src/engines/fabricx/}),
  fonttitle=\bfseries, arc=4pt]
\begin{itemize}[leftmargin=*]
\item \textbf{Crypto material is copied, not generated.} The current engine resolves a crypto source directory from \texttt{FABLO\_FABRICX\_CRYPTO\_SOURCE} or known relative sample paths and copies it into the target directory. My mentorship plan replaces this with a staged crypto-generation pipeline based on \texttt{cryptogen}, Fabric CA enrollment, and \texttt{tokengen}.
\item \textbf{Only xdev bundled topology is supported.} The generated compose file starts one bundled \texttt{committer} service and per-node config files. This is enough for a local developer workflow, but it does not yet model the decomposed orderer and committer family visible in the Ansible inventory. The post-MVP plan extends the schema with explicit topology modes. The MVP is scoped to the bundled \texttt{xdev} topology (single container with db, orderer, committer), not decomposed multi-service setups, to ensure a mergeable first implementation.
\item \textbf{Orderer identity defaults are sample-based.} The compose template currently uses \texttt{FABLO\_FABRICX\_ORDERER\_MSP\_ID} and \texttt{FABLO\_FABRICX\_ORDERER\_MSP\_DIR\_REL} with defaults derived from a sample admin MSP path. My longer-term fix is to promote orderer identity and MSP ownership into explicit configuration and generated crypto outputs.
\item \textbf{Runtime compatibility is image-driven, not profile-driven.} The current schema carries \texttt{global.fabricVersion}, but Fabric-X behavior is effectively governed by the infrastructure image and supporting crypto and tool versions. My next step is to introduce explicit compatibility profiles and a machine-readable matrix.
\item \textbf{No Fabric-X-specific reset or namespace tooling is merged yet.} The MVP focuses on \texttt{generate}, \texttt{up}, \texttt{down}, and \texttt{status}. Reset and future \texttt{fxconfig}-style helpers should be added only after the local bootstrap path is stable.
\end{itemize}
None of these limitations affect classic Fabric behavior.
\end{tcolorbox}

The important point is not that the MVP has limitations; all credible MVPs do. The important point is that each limitation is already isolated and explained. The engine split ensures that classic Fabric behavior remains unaffected. I have also defined the replacement path for each shortcut directly in this proposal, which keeps the mentorship plan reviewable and lowers the risk of architectural drift.

\section{Pre-Application Research}

\begin{tcolorbox}[colback=green!5!white, colframe=green!60!black,
  title=Pre-Application Research Completed, fonttitle=\bfseries, arc=4pt]
Before writing this proposal and building the proof of concept, the
following research was completed:
\begin{itemize}[leftmargin=*]
  \item Read the Fablo command and generation path across \texttt{src/commands/generate/index.ts}, \texttt{src/commands/up/index.ts}, \texttt{src/commands/down/index.ts}, \texttt{src/commands/init/index.ts}, \texttt{src/engines/engine.ts}, \texttt{src/engines/engine-loader.ts}, \texttt{src/engines/classic/index.ts}, and \texttt{src/setup-docker/index.ts}; confirmed that the CLI can delegate cleanly to a pluggable engine.
  \item Read the current Fabric-X engine implementation and templates in \texttt{src/engines/fabricx/index.ts}, \texttt{src/engines/fabricx/schema/fabricx-schema-v1.json}, \texttt{docker-compose.xdev.yaml.ejs}, \texttt{routing-config.yaml.ejs}, and \texttt{core.yaml.ejs}; confirmed that a real schema-driven generation flow already exists.
  \item Read Fablo's test and feature surface through \texttt{SUPPORTED\_FEATURES.md}, \texttt{e2e-network/docker/test-01-v3-simple.sh}, and \texttt{e2e-network/fabricx/test-01-fabricx-simple.sh}; confirmed that Fabric-X can and should mirror the existing integration-test style.
  \item Read Fabric-X architecture and deployment material through \texttt{fabric-x-repo/README.md}, \texttt{samples/tokens/compose-xdev.yml}, \texttt{samples/tokens/ansible/inventory/fabric-x.yaml}, \texttt{samples/tokens/scripts/gen\_crypto.sh}, \texttt{tools/fxconfig/}, and \texttt{go.mod}; confirmed the gap between bundled local topology and decomposed long-term topology.
  \item Identified four key architectural questions that I then addressed directly in my design work: crypto generation, orderer identity management, topology expansion, operational configuration surface, and compatibility profile design.
  \item Built and hardened a working proof of concept: engine-driven Fabric-X support, \texttt{fablo init --fabricx}, shared target-dir resolution, explicit generate-before-up behavior, precondition checks, port defaults, TLS-aware template rendering, improved polling, stricter types and schema validation, and a Fabric-X e2e test.
  \item Produced diagrams, implementation notes, and a running history of the changes so I could validate command flow, architecture, implementation history, and cumulative design decisions before drafting the formal proposal.
\end{itemize}
\end{tcolorbox}

This pre-application work is important because it changes the proposal from speculative to evidence-based. The working PoC is not presented as final software, but it is strong proof that the architectural direction is sound and that the candidate has already navigated the difficult parts of the codebase.

\section{Why I Am the Right Candidate}

\begin{table}[h]
\centering
\begin{tabularx}{\textwidth}{|p{3.2cm}|X|}
\hline
\textbf{Area} & \textbf{Relevant experience} \\
\hline
TypeScript / Node.js & Built full-stack systems in JavaScript and TypeScript, including Campus Marketplace and AuthrNet; already modified Fablo's TypeScript command, engine, and template-generation code. \\
\hline
Go & Read and reasoned through Fabric-X's Go-based \texttt{fxconfig} tooling, \texttt{go.mod}, and integration tests; experienced enough with Go ecosystems to understand versioning, module boundaries, and CLI tool integration. \\
\hline
Bash / CLI tooling & Comfortable with shell-driven workflows through e2e test reading and implementation work; the Fabric-X proof of concept already extends Fablo CLI behavior and its Bash-based test path. \\
\hline
Docker / compose & Extensive Docker and GitHub Actions use in prior projects; directly relevant to compose generation, multi-service orchestration, and startup-debugging for Fabric-X networks. \\
\hline
Distributed systems & Worked on decentralized and backend-heavy projects; understands multi-service startup ordering, identity boundaries, and infrastructure debugging, which are central to Fabric-X. \\
\hline
Open source & More than 800 GitHub contributions, FOSS mentoring, and active collaborative development experience; relevant to delivering reviewable, maintainable work in an upstream community. \\
\hline
Debugging multi-service systems & Experience integrating APIs, services, authentication, databases, and deployment tooling; the current Fabric-X work already involved tracing runtime assumptions across generated config, Docker Compose, and sample crypto layout. \\
\hline
\end{tabularx}
\caption{Skills and experience mapping to project requirements}
\label{tab:skills}
\end{table}

My fit for this project comes from the overlap between implementation skills and evidence of already having done the architectural homework. The Fabric-X plus Fablo integration is not a narrow blockchain task. It requires TypeScript CLI work, schema design, template generation, Docker runtime understanding, test automation, and careful reading of a second codebase whose core tooling is in Go. That blend matches my current profile well. My background in TypeScript, JavaScript, backend systems, and Docker-based workflows means I am comfortable making changes across the exact layers this project touches.

The strongest evidence is the work already present in my proof of concept and in the implementation described throughout this proposal. I did not stop at reading READMEs and writing a high-level plan. I implemented the engine split cleanup tasks, updated command routing, hardened the Fabric-X generator, fixed preconditions, aligned TLS handling, tightened schema and type boundaries, improved test coverage, and worked through the open design questions that still matter. The architectural diagrams I created were not decorative; they were part of understanding how the current implementation behaves and where its seams are.

My project history also maps cleanly onto the needs of this mentorship. Campus Marketplace involved full-stack system design, authentication, RBAC, and hybrid data-layer work, which is relevant because Fablo integration work is fundamentally about turning a user-facing developer surface into correct infrastructure behavior. AuthrNet is relevant because it required blockchain-oriented reasoning, identity flows, and coordinating a user-facing application with decentralized primitives. My internship and hackathon backend work are relevant because they demonstrate a pattern of working across interfaces, services, and deployment environments instead of remaining inside a single feature layer.

Open source fit matters just as much as raw technical fit. This project needs someone who can separate what should be merged now from what should be deferred, document known limitations honestly, and structure work so maintainers can review it in small, coherent pieces. The existing proof of concept, research notes, and documentation updates show that I approach the work that way. I am already treating the Fabric-X path as something that must preserve classic Fabric behavior, maintain a clean extension seam, and stay reviewable by upstream maintainers. That mindset is as important as writing code quickly.

\section{Risk Register}

\begin{table}[h]
\centering
\begin{tabularx}{\textwidth}{|p{3cm}|X|p{2cm}|}
\hline
\textbf{Risk} & \textbf{Mitigation strategy} & \textbf{Probability} \\
\hline
Fabric-X API surfaces change before stable release & Keep version-sensitive behavior isolated inside \texttt{src/engines/fabricx/}; pin sample images explicitly; document a compatibility profile strategy so schema contracts can remain stable across runtime image updates. & Medium \\
\hline
Crypto generation for FSC nodes is underdocumented & Treat \texttt{samples/tokens/scripts/gen\_crypto.sh} as the executable reference, implement the crypto pipeline in stages, and keep a fallback path for externally supplied crypto material during transition. & High \\
\hline
xdev bundled topology differs from final decomposed topology & Declare \texttt{xdev} as the MVP scope, design the schema with explicit topology evolution in mind, and add decomposed support only after the bundled engine path is merged and tested. & Medium \\
\hline
Docker networking for decomposed services is complex & Start from the known Ansible inventory and bundled sample port map; introduce decomposed services incrementally rather than as a single large compose rewrite. & Low \\
\hline
Stretch goals crowd core work & Core must be merged before Week 13 stretch begins; use the engine boundary and documented limitation list to keep scope decisions explicit. & Medium (mitigated) \\
\hline
\end{tabularx}
\caption{Risk register with mitigation strategies}
\label{tab:risks}
\end{table}

The most important feature of this risk register is that I tie every item to something directly observed in the repositories rather than to generic software-project uncertainty. Version drift is visible already between the sample image lineage and current Go module state. Crypto complexity is visible in the sample shell script. Topology mismatch is visible between \texttt{compose-xdev.yml} and the Ansible inventory. Those are real risks, which makes them manageable.

\section{Success Metrics}

\begin{tcolorbox}[colback=green!5!white, colframe=green!60!black,
  title=Core Deliverables --- all must be complete, fonttitle=\bfseries, arc=4pt]
\begin{itemize}[leftmargin=*]
  \item[\checkmark] Design proposal document accepted and merged
  \item[\checkmark] \texttt{fablo init --fabricx} generates valid, schema-identified config
  \item[\checkmark] \texttt{fablo generate} produces complete docker-compose, routing-config, and per-node core.yaml
  \item[\checkmark] \texttt{fablo up / down / status} lifecycle working end-to-end
  \item[\checkmark] E2E test suite passing, following existing Fablo test structure
  \item[\checkmark] Contributor and user documentation merged
  \item[\checkmark] At least two community members have successfully run a Fabric-X network using the implementation
\end{itemize}
\end{tcolorbox}

\begin{tcolorbox}[colback=blue!5!white, colframe=blue!60!black,
  title=Stretch Goals --- conditional on core completion by Week 12,
  fonttitle=\bfseries, arc=4pt]
\begin{itemize}[leftmargin=*]
  \item Multi-org Fabric-X topology support
  \item Separate-container committer, validator, verifier, coordinator, and orderer deployment
  \item CI integration example for Fabric-X networks
  \item Fablo UI support for Fabric-X configs and topology visualization
\end{itemize}
\end{tcolorbox}
colorbox}
